\section{Conclusion}\label{sec:conclusion}

\kit{} gives a team with one GPU three things: a runner that has been attacked and whose misses are published, checks
that compare a local serving stack with the scorer down to token ids, and a pass rule with known error rates. This
extended version adds what the kit was used for: \RunsTotal{} agent runs in \BlocksTotal{} blocks and \ArmsTotal{}
arms, and the \CaseRows{} rows of the case study.

Most of what the record shows is the absence of a difference. Six configurations of a 31B agent lay between
\PubMinSeventy{} and \PubMaxSeventy{} of 70 development tasks, and no pairwise difference was detected. Four
prompt divergences changed the E4B solve rate by no detectable amount. Three changed E4B agents led their screens and
were not confirmed on new tasks, with gains of $\GainOne$, $\GainTwo$ and $\GainThree$ points. The third of them is
the closest case: its gain is large against run-to-run noise and does not pass the test across repositories. We do
not call it an improvement.

The record also explains why this is the expected outcome at this scale. On a fixed set of tasks the small model is
nearly deterministic at the level of a task: about two thirds of the tasks were solved in no run, and the baseline
repeated its total within \RepeatHDiff{} points. A screen of 96 runs per arm cannot rank variants, and its leader is
ahead by selection. A hosted score is one draw with a standard deviation near \HostedSdTasks{} tasks of 58. The
instrument that can decide is a confirmation on tasks from repositories that no earlier experiment has used, and
each one spends its tasks. The validated pool that we used is almost spent.

We release every run-level result so that others can recompute these numbers, plan the size of their own
comparisons from the measured noise, and avoid paying again for the same screens.
